\subsection*{MT-M26. Trace class on the specified strip Hilbert quotient}
Fix one primitive unitary Hecke function and the notation of
MT-M18 and MT-M24. Let $s\ge0$, $a\in\mathbf R$, and let
$I=[u,v]$ be a compact interval with $u<v$. Define
\[
 \begin{split}
 w_{I,s}(x)&=\max\{e^{2ux},e^{2vx}\}(1+x^2)^{s/2},\\
 H_{I,s}&=L^2(\mathbf R,w_{I,s}(x)\,dx),\\
 J_{a,I,s}&=\overline{T_a\mathcal I_\eta}^{\,H_{I,s}},
 \qquad Z_{a,I,s}=H_{I,s}/J_{a,I,s},
 \qquad T_af(x)=e^{ax}f(x).
 \end{split}                                                    \tag{M26.1}
\]
Assume there is a closed strip $[c_0,c_1]$ containing the real
parts of all zeros of $\Xi_\eta$, with
\[
                 [c_0-a,c_1-a]\subset(u,v).
                                                               \tag{M26.2}
\]
Then the natural map $Q_\eta\to Z_{a,I,s}$ is continuous,
injective and has dense image. More precisely,
\[
 J_{a,I,s}
 =\{h\in H_{I,s}:(\mathcal Mh)^{(k)}(\rho-a)=0
       \text{ for all zeros }\rho\text{ and }0\le k<m_\rho\}.
                                                               \tag{M26.3}
\]
The Mellin transforms and derivatives in (M26.3) are defined
by absolutely convergent integrals inside the strip $I$.
For every $g\in\mathcal E$, convolution induces a trace-class
operator on this \emph{actual} Hilbert quotient, and
\[
       \operatorname{Tr}_{Z_{a,I,s}}C_g
          =\sum_{\Xi_\eta(\rho)=0}
                        m_\rho(\mathcal Mg)(\rho-a),
                                                               \tag{M26.4}
\]
with absolute convergence. The trace-class estimates are
uniform when $g$ ranges over a bounded subset of $\mathcal E$.

Unconditionally one may take $[c_0,c_1]=[0,1]$, so the result
holds whenever $[0,1]\subset a+\operatorname{int}I$.
Under GRH for this one function, one may take $c_0=c_1=1/2$.
In particular, under the conditional input $M$, the result
holds at $a=1/2$ for every $I=[-\alpha,\alpha]$ with
$\alpha>0$, for each fixed Hecke character separately.

\begin{proof}
We first specify the adapted blocks needed inside the available
strip. Choose fixed real numbers
\[
        a+u<b_-<c_0\le c_1<b_+<a+v.
\]
Use the selected horizontal heights $h_n$ from MT-M18 and
rectangles $\Omega_n=(b_-,b_+)\times(h_n,h_{n+1})$.
The proof of MT-M18 gives the same block theorem for these
vertical edges, as the following estimates make explicit.
The horizontal minimum-modulus estimate there was proved
on every fixed horizontal segment, so it applies on
$[b_-,b_+]$. On either vertical edge the distance from every
zero is at least
$\delta=\min(c_0-b_-,b_+-c_1)>0$.
In the Blaschke-product proof of MT-M18, each factor on these
edges is therefore at least
$\min(1,\delta)/(2R)$, where $R$ is the fixed disk radius
used there. There are $O(\log(2+|n|))$ such factors.
Their product has a polynomial lower bound. The same
zero-free factor estimate and the two-sided Stirling bound
then give
\[
 \inf_{w\in\partial\Omega_n}|\Xi_\eta(w)|
    \ge c(2+|n|)^{-B}e^{-\kappa|n|}.
                                                               \tag{M26.5}
\]
At bounded height the edges contain no zeros and compactness
supplies the bound. The fixed width and bounded height of the
rectangles give uniformly bounded perimeter. The entire
difference-kernel construction in MT-M18, its Gaussian estimate,
and its quotient-infimum argument consequently give finite
Banach algebras $J_n$ and lifts $\mathcal T_n:J_n\to\mathcal E$
with
\[
 \begin{gathered}
 d_n=\dim J_n=O(\log(2+|n|)),\qquad
 R_m\mathcal T_n=\begin{cases}\operatorname{id}_{J_n},&m=n,\\0,&m\ne n,
                         \end{cases}\\
 p_{N,j}(\mathcal T_n z)
       \le C_{N,j}(2+|n|)^{B_{N,j}}\|z\|_{J_n}.
 \end{gathered}                                                 \tag{M26.6}
\]
Here $J_n$ is the quotient of the holomorphic supremum-norm
algebra on the closed rectangle by its full block jets.
The absolute synthesis, closed-ideal equality and two-sided
block isomorphism of MT-M18 also hold: their proof uses
(M26.5), bounded perimeter and inclusion of every zero in
exactly one rectangle, all just verified. In particular the
lifts are not unstable individual Hermite interpolation
vectors. This establishes the precise version of MT-M18
being used below.

For $h\in H_{I,s}$ its Mellin integral is holomorphic on
$u<\Re z<v$. If $\sigma$ ranges over a compact subinterval
of $(u,v)$, then for every integer $k\ge0$,
\[
 \begin{split}
 |(\mathcal Mh)^{(k)}(\sigma+it)|
 &\le\|h\|_{H_{I,s}}
   \left(\int_{\mathbf R}
     \frac{|x|^{2k}e^{2\sigma x}}{w_{I,s}(x)}\,dx\right)^{1/2}\\
 &\le C_k\|h\|_{H_{I,s}},
 \end{split}                                                    \tag{M26.7}
\]
uniformly in $t$. The integral converges by the positive
distance to both endpoints, uniformly in that compact
subinterval. The same estimate justifies differentiation
under the integral on compact subsets of the strip.
Consequently
\[
 R_n^H(h)=\bigl[w\longmapsto(\mathcal Mh)(w-a)\bigr]\in J_n
 \quad\hbox{satisfies}\quad
       \|R_n^H(h)\|_{J_n}\le C\|h\|_{H_{I,s}},
                                                               \tag{M26.8}
\]
where $C$ is independent of $n$. This follows by bounding
the quotient norm by the supremum of this representative
on the contour. For $h=T_af$ the representative is
$\mathcal Mf(w)$, so $R_n^H$ vanishes on $T_a\mathcal I_\eta$,
and by continuity on its Hilbert closure. It induces a
bounded map $r_n:Z_{a,I,s}\to J_n$ with norm at most $C$.

Let $\pi:H_{I,s}\to Z_{a,I,s}$ be the quotient map, and put
$s_n=\pi T_a\mathcal T_n$. The map $T_a:\mathcal E\to H_{I,s}$
is continuous: a seminorm $p_{N,0}$ with
$N>|a|+\max(|u|,|v|)+1$ bounds its norm by an integrable
exponential times the fixed polynomial weight. Thus (M26.6)
implies, for some fixed constants $C,B$,
\[
 \|s_n\|_{J_n\to Z_{a,I,s}}\le C(2+|n|)^B,
 \qquad r_m s_n=\delta_{mn}\operatorname{id}_{J_n}.
                                                               \tag{M26.9}
\]
The last identity follows directly from the jets of the
lifts. In particular each block is actually realized in
the Hilbert quotient.

Convolution by $g\in\mathcal E$ is bounded on $H_{I,s}$.
Indeed, with $B_I=\max(|u|,|v|)$, the elementary polynomial
weight inequality used in MT-M25 gives
\[
 \|h(\,\cdot-y)\|_{H_{I,s}}
 \le2^{s/4}e^{B_I|y|}(1+y^2)^{s/4}\|h\|_{H_{I,s}}.
                                                               \tag{M26.10}
\]
Minkowski's inequality integrates this bound against
$|g(y)|$, and the resulting integral is finite.
The subspace $T_a\mathcal I_\eta$ is invariant under this
convolution because
\[
           g*T_af=T_a\bigl((e^{-a\cdot}g)*f\bigr)
                                                               \tag{M26.11}
\]
and $\mathcal I_\eta$ is a convolution ideal in $\mathcal E$.
Its Hilbert closure is therefore invariant as well. Hence
$C_g$ induces a bounded operator on $Z_{a,I,s}$.

Let $M_{g,n}$ be multiplication in $J_n$ by the block class
of $w\mapsto\mathcal Mg(w-a)$. Submultiplicativity of the
quotient algebra norm shows
\[
 \|M_{g,n}\|\le b_n(g):=
       \sup_{w\in\partial\Omega_n}|\mathcal Mg(w-a)|.
                                                               \tag{M26.12}
\]
For each $A>0$, $b_n(g)\le C_A(g)(2+|n|)^{-A}$,
with $C_A(g)$ bounded on bounded subsets of $\mathcal E$,
by its strip-Schwartz Mellin transform. These constants
are bounded by finitely many $\mathcal E$ seminorms:
integrating by parts sufficiently many times on the
fixed strip bounds the high ordinates, and direct
integration bounds the remaining compact set. The finite-rank
operators $s_nM_{g,n}r_n$ have rank at most $d_n$ and
satisfy the Hilbert trace-norm bound
\[
 \|s_nM_{g,n}r_n\|_1
       \le C d_n(2+|n|)^B b_n(g).
                                                               \tag{M26.13}
\]
Here we used the elementary inequality
$\|T\|_1\le\operatorname{rank}(T)\|T\|$ for a finite-rank
operator on a Hilbert space. Thus the series
\[
                     \sum_n s_nM_{g,n}r_n
                                                               \tag{M26.14}
\]
converges in trace norm, uniformly for $g$ in bounded test
sets. These same estimates show that
$g\mapsto C_g$ is continuous into the trace-class norm.

It equals the actual induced convolution. To verify this
without assuming spectral completeness in the Hilbert
space, first take a vector $\pi T_af$ with $f\in\mathcal E$.
Equation (M26.11) and the absolute block synthesis in the
\emph{already proved smooth quotient} give the identity
with (M26.14) on this vector. Applying $\pi T_a$ is
legitimate because it is continuous from that smooth
quotient into the Hilbert quotient; no injectivity has
been assumed. These vectors are dense: $T_a\mathcal E$
contains $C_c^\infty$, which is dense in the weighted
$L^2$ space by truncation and compact approximation.
Both operators are bounded, so the identity holds on
all of $Z_{a,I,s}$.

We can now prove the exact closure statement (M26.3).
One inclusion follows already from (M26.7). For the other,
if every zero jet of $h\in H_{I,s}$ vanishes, then
$r_n[ h ]=0$ for every $n$. The operator identity just
proved gives $C_g[h]=0$ for every $g\in\mathcal E$.
Take a smooth compactly supported approximate identity
$g_\epsilon$ of integral one. Translation is strongly
continuous on $H_{I,s}$: this holds first for compactly
supported smooth functions, then for all functions by
their density and (M26.10). That estimate is uniformly
bounded on small $y$, so
$g_\epsilon*h\to h$ in $H_{I,s}$.
Passing to the quotient gives $[h]=0$, proving (M26.3).
In particular $Q_\eta\to Z_{a,I,s}$ is injective, since
all zero jets of a smooth vector still determine its
class by MT-M18. Continuity and density were proved above.

Finally the finite-rank trace identity and (M26.9) give
\[
 \operatorname{Tr}(s_nM_{g,n}r_n)
       =\operatorname{Tr}_{J_n}(M_{g,n}r_ns_n)
       =\operatorname{Tr}_{J_n}M_{g,n}.
\]
In the local algebra at a zero $\rho$ of multiplicity
$m_\rho$, multiplication by a holomorphic function is
triangular in the powers of $w-\rho$, with that function's
value at $\rho$ repeated $m_\rho$ times on the diagonal.
This computes the trace as
$\sum_{\rho\in\Omega_n}m_\rho\mathcal Mg(\rho-a)$.
Summing is legitimate by (M26.13), and the scalar sum
is also absolutely convergent since $d_n=O(\log(2+|n|))$
and the Mellin transform is rapidly decreasing on the
fixed strip. This proves (M26.4).
\end{proof}

The injection from the smooth quotient has proper dense
image. Indeed $\Xi_\eta$ has infinitely many nontrivial
zeros, hence infinitely many nonzero $J_n$. The rapid
block space of MT-M18 is not normable: if one norm gave
its topology, it would be bounded above by a constant
times some $q_K$, while $q_{K+1}$ would be bounded by
a constant times that norm. A unit vector in a single
block at height $n$, scaled by $(2+|n|)^{-K}$, contradicts
these two inequalities as $|n|\to\infty$. Were the
continuous injection onto the Hilbert quotient surjective,
the Fr\'echet open mapping theorem would make the smooth
quotient normable. Thus this is a faithful Hilbert
completion, not an identification of its norm topology
with the stronger smooth topology.

At $a=1/2$ put $h^*(x)=\overline{h(-x)}$. Under GRH,
(M26.4) yields
\[
 \operatorname{Tr}C_{h*h^*}
      =\sum_\rho m_\rho|\mathcal Mh(\rho-1/2)|^2\ge0.
                                                               \tag{M26.15}
\]
On any strip satisfying the unconditional sufficient
condition $[0,1]\subset1/2+\operatorname{int}I$, positivity
of these traces for every $h\in\mathcal E$ is equivalent
to GRH for the fixed function. For if
$\rho^\dagger=1-\overline\rho\ne\rho$ is an off-line pair,
the unitary functional equation gives equal multiplicities.
MT-M18 finite-block interpolation, followed by $T_{1/2}$,
produces $h$ with Mellin values $1,-1$ at
$\rho-1/2,\rho^\dagger-1/2$, respectively, and zero at
all other zeros. Formula (M26.4) then gives
$\operatorname{Tr}C_{h*h^*}=-2m_\rho$.
No positivity of the operator $C_{h*h^*}$ itself is claimed;
the faithful full-jet obstruction of MT-M20 still applies.

\paragraph{Literal source interface and remaining scope.}
Masullo's norm (9.2.2.31), original PDF p.~416, becomes
(M26.1) after $x=\log|g|$ in the fixed compact-character
sector; the supremum over subintervals is attained at
the endpoints of the compact interval $I$. His closure
Definition~10.1.1.1, p.~434, becomes $J_{a,I,s}$ by
MT-M24. For $I=[-\alpha,\alpha]$ and $a>0$, his stated
parameter range, the resulting space
is his specified quotient (10.4.1.3), original PDF p.~456,
in that sector. The convolution convention is exactly
(M24.11); an idele-class test uses the scalar function
$k_\beta$, so its trace is evaluated at the correspondingly
shifted quasi-character. Compactly supported smooth tests
are included in $\mathcal E$.

This independently supplies fixed-character spectral
trace class on those literal quotients; it does not invoke
Theorem~10.3.2.1 or the draft's general Schauder claims.
All constants may depend on the fixed field and character.
A simultaneous summation over varying compact characters
would require additional uniform conductor and archimedean
parameter estimates; that is not inferred from a
fixed-character result. Neither the trace formula nor
the space identification proves $M$. The remaining
positivity input is still conditional.

In the inherited cote~1 / SA23 geo-H-lidskii and geo-H-genus
crosswalk, the map from the completed-Hecke spectral
divisor to a specified adelic Hilbert quotient is now
explicit and faithful whenever the containing strip has
interior room. The line quotient MT-M25 explains why
that interior condition cannot simply be discarded.
